\documentclass[11pt,a4paper]{article}
\usepackage[margin=24mm]{geometry}
\usepackage{amsmath,amssymb,amsthm,hyperref}
\newtheorem{theorem}{Theorem}\newtheorem{lemma}[theorem]{Lemma}\newtheorem{proposition}[theorem]{Proposition}
\newcommand{\gen}[1]{\langle #1\rangle}\newcommand{\im}{\operatorname{im}}\newcommand{\GL}{\operatorname{GL}}\newcommand{\PGL}{\operatorname{PGL}}
\title{Projective involutions and support spaces in small dimension}
\author{Human Compression project: recovery of L011}\date{}
\begin{document}\maketitle
This paper gives the geometric exclusions used in the simple-group classification
of Leemans and Vandenschrick~\cite{LV}. The dimension-three arguments rework
Brooksbank and Vicinsky~\cite{BV}, including the projective qualification. The
six-dimensional argument reworks Brooksbank and Leemans~\cite{BL} for both
quadratic types. Attribution to these arguments is retained. The deep
rank-three existence classification is a separate input of the application.

A string marking is a sequence of distinct nonidentity involutions generating
a group, with generators at distance at least two commuting. A string
\emph{C-group} marking also has the intersection property. Every subsequence
of a C-group marking is independent: none of its generators lies in the
subgroup generated by the others. We use this implication explicitly.

\section{Projective dimension three}
Let $F$ be a finite field and $V=F^3$. A projective involution has a canonical
linear involution representative. Indeed, if $X^2=\mu I$, then
$(\det X)^2=\mu^3$, so $\mu$ is a square. In odd characteristic normalize
$X$ so its two-dimensional eigenspace has eigenvalue $1$; the other eigenvalue
is $-1$. In characteristic two normalize its unique eigenvalue to $1$.
Then $X=I+N$ with $N^2=0$ and $\operatorname{rank}N=1$.
In either characteristic the representative is uniquely determined by the
projectivity and the specified normalization. Conjugation preserves it.
Consequently commuting projectivities have commuting canonical representatives.

In odd characteristic write $\delta_X$ for the one-dimensional eigenspace
and $\Delta_X$ for the other eigenspace. For distinct involutory projectivities,
commutation is equivalent to
\[
 \delta_X\subset\Delta_Y,\qquad \delta_Y\subset\Delta_X.
\tag{1}
\]
This follows by preserving eigenspaces; their unequal dimensions prevent an
exchange, and equal centers would force equal axes when the representatives
commute. The word ``distinct'' matters: the criterion does not apply to an
involution compared with itself.
In characteristic two let $\delta_X=\im N$ and $\Delta_X=\ker N$.
Here $\delta_X\subset\Delta_X$, and two transvections commute precisely when
their centers coincide or their axes coincide. To check this, write
$N=u f$, $M=v h$. If $u,v$ are independent, $NM=MN$ forces both products
to vanish, so both axes contain $\gen{u,v}$ and coincide. Equal centers
conversely give zero products, as do equal axes.

\begin{lemma}[Connectedness and absolute irreducibility]\label{connected}
If a projective group on $F^3$ is irreducible and is generated by a string
marking, then its canonical linear lift is absolutely irreducible and every
adjacent pair is noncommuting.
\end{lemma}
\begin{proof}
The lift is irreducible because invariant linear subspaces are the same as
invariant projective subspaces. Its commuting endomorphism algebra is a finite
division algebra. Any commuting endomorphism preserves the one-dimensional
space $\delta_X$ of one of the canonical involutions. Restriction to that line
is injective: a nonzero commuting endomorphism is invertible. The commuting
algebra is therefore exactly $F$. The standard finite-dimensional algebra
argument now gives absolute irreducibility: the faithfully acting algebra is
semisimple (its radical annihilates a simple module), has one simple factor,
and is a matrix algebra over its commuting division field, here $F$.
Thus its image is the full endomorphism algebra and remains so after extending
scalars.

If an adjacent pair commuted, the string would split into two nonempty parts
whose generated linear groups commute. Over an algebraic closure an irreducible
module of their direct product is a tensor product of irreducible modules.
For completeness this is the matrix-algebra decomposition of the two commuting
simple algebra images; each is the full matrix algebra on its tensor factor.
The product of the factor dimensions is three, so one factor is one-dimensional.
That group then acts by scalars, contrary to its containing a nonidentity
projective generator. The same argument applies to any cut.
\end{proof}

\begin{theorem}[Projective form obstruction]\label{three}
An irreducible string C-group in $\PGL_3(F)$ has rank at most four,
$F$ has odd characteristic, and its projective action lifts to a group
preserving a nondegenerate symmetric bilinear form on $F^3$.
\end{theorem}
\begin{proof}
Write the canonical lifts as $T_0,\ldots,T_{r-1}$. Rank one is reducible.
For rank two the two axes intersect in a nonzero invariant line. We may thus
assume $r\ge3$ and use Lemma~\ref{connected}.

First suppose the characteristic is two. The centers and axes of $T_0,T_1$
are distinct, since those involutions do not commute. If neither center lies
in the other's axis, no transvection can commute with both: it would have
to share the center or axis of each, which would force one of those
incidences. Thus $r\le3$. The commuting endpoints of a triple share a center
or an axis. In the former case their center together with the middle center
spans an invariant plane; in the latter case their common axis intersects
the middle axis in a common fixed line. In both cases the triple is reducible.

In the remaining case, duality permits us to assume
$\delta_0\subset\Delta_1$. The opposite incidence is impossible, since it
would make both distinct axes the plane spanned by the centers.
Every transvection commuting with both $T_0,T_1$ then has center
$\delta_0$ and axis $\Delta_1$. These transvections commute with each other,
so noncommuting adjacency gives $r\le4$. Also $T_2$ shares the center or
axis of $T_0$ and hence fixes $\delta_0$ pointwise. Every later generator
has center $\delta_0$ and fixes it, as do $T_0,T_1$. This contradicts
irreducibility. Characteristic two is excluded.

Now suppose the characteristic is odd. If the first two centers coincide,
all other generators preserve that center, by commuting with $T_0$ or $T_1$.
If their axes coincide, the same argument gives a common plane. Hence both
centers and axes are distinct. By (1), a projective involution commuting
with both is uniquely determined, if it exists: its center is
$\Delta_0\cap\Delta_1$ and its axis is $\gen{\delta_0,\delta_1}$.
Two generators with indices at least three would therefore coincide. Independence gives
$r\le4$.

For rank three, simultaneously diagonalize the commuting endpoints, giving
their exceptional eigenspaces the first and third coordinate lines.
The space of diagonal symmetric forms has dimension three. The symmetric
forms preserved by $T_1$ have dimension four: in its eigenbasis they have
one block of size one and one symmetric block of size two. Both spaces lie
in the six-dimensional symmetric-matrix space, so their intersection is
nonzero. Its radical is invariant under the full group and therefore zero.

For rank four, simultaneously diagonalize $T_0,T_3$ in the same way.
The matrix $T_1$ preserves the coordinate plane $\gen{e_1,e_2}$ and line
$\gen{e_3}$; $T_2$ preserves $\gen{e_1}$ and $\gen{e_2,e_3}$.
Consider a diagonal symmetric form $B=\operatorname{diag}(x,y,z)$.
For an involution $T$, invariance is equivalent to $BT=T^{\mathsf T}B$.
The two block matrices impose at most one homogeneous equation each:
\[
 (T_1)_{12}x=(T_1)_{21}y,\qquad
 (T_2)_{23}y=(T_2)_{32}z.
\]
There is a nonzero solution over $F$. Its radical is again invariant under
all four generators, hence zero. This argument does not presume that the
first rank-three subgroup is irreducible, or that the initial dihedral
subgroup acts semisimply.
\end{proof}

\begin{proposition}\label{psl3}
No nonabelian simple group $\operatorname{PSL}_3(q)$ or
$\operatorname{PSU}_3(q)$ has a string C-group marking of any rank.
\end{proposition}
\begin{proof}
Their natural projective actions over $\mathbb F_q$ and $\mathbb F_{q^2}$
are irreducible, a standard natural-module fact. Even characteristic is
excluded by Theorem~\ref{three}. For odd $q\ge3$, a lifting group as in that
theorem would lie in the orthogonal \emph{isometry} group, whose order in
dimension three over $\mathbb F_t$ is $2t(t^2-1)$.
The usual classical-group orders give
\[
 \frac{|O_3(q)|}{|\operatorname{PSL}_3(q)|}
 =\frac{2\gcd(3,q-1)}{q^2(q^3-1)}<1,
\]
and
\[
 \frac{|O_3(q^2)|}{|\operatorname{PSU}_3(q)|}
 =\frac{2\gcd(3,q+1)(q^2+1)}{q(q^3+1)}
 \le\frac{6(q^2+1)}{q(q^3+1)}<1.
\]
A group projecting onto the indicated simple group cannot have smaller order.
The natural-module and order formulas are retained classical inputs;
no finite-field sampling is used.
\end{proof}

\section{Orthogonal dimension six in characteristic two}
Let $(V,Q)$ be a nonsingular six-dimensional quadratic space over the finite
field $F$ of characteristic two, of either type. Write $B$ for its polar
alternating form. We retain the classical Dickson-invariant identity
\[
 \Omega(V)=\{g\in O(V):\operatorname{rank}(g-I)\text{ is even}\}.
\tag{2}
\]
Together with the two Klein identifications used below, this is the classical
orthogonal-group input. The remaining support calculation is given explicitly.

\begin{lemma}[Involution supports]\label{support}
For a nonidentity involution $g\in\Omega(V)$ its support
$U=\im(g-I)$ has dimension two, is totally isotropic for $B$, and
$\ker(g-I)=U^\perp$. If two such involutions commute, their supports intersect
nontrivially. If they do not commute and their supports intersect, their
intersection is a nonsingular line for $Q$.
\end{lemma}
\begin{proof}
Put $N=g-I$. Since $N^2=0$, its rank is at most three; (2) makes it two.
Because $g$ is an involutory isometry, $N$ is self-adjoint for $B$.
It follows that $\ker N=U^\perp$ and that $B(Nx,Ny)=B(x,N^2y)=0$.
Choose a basis $e,b$ of $U$ with $Q(e)=0$; such an $e$ exists since the
field is perfect and the polar form on $U$ is zero. Self-adjointness and
$Q(x+Nx)=Q(x)$ show that, for some $\alpha\ne0$,
\[
 Nx=\alpha B(x,b)e+\alpha B(x,e)b
       +\alpha^2Q(b)B(x,e)e.\tag{3}
\]
Here is the coefficient verification. Write a self-adjoint map with image
in $U$ as
$Nx=\alpha B(x,b)e+\alpha B(x,e)b+\beta B(x,e)e+\gamma B(x,b)b$.
The quadratic identity gives $\beta=\alpha^2Q(b)$ and
$\gamma=\gamma^2Q(b)$. If $\gamma\ne0$, the coefficient matrix has
determinant $\alpha^2(1+\gamma Q(b))=0$, contrary to rank two.
Thus $\gamma=0$, and rank two forces $\alpha\ne0$.

For commuting $N,M$ with disjoint images, $NM=MN$ has image in their
intersection and is zero. Their two support planes are then mutually
orthogonal, so their sum is a four-dimensional totally isotropic space for
the nonsingular alternating form $B$ on $V$, impossible. Thus their supports
intersect. If the common line is singular, take its vector as $e$ in (3)
for both maps. Substitution gives $NM=MN$, both scalar multiples of the map
$x\mapsto B(x,e)e$. Finally, equal support planes also give commuting maps,
since each map kills that plane. Hence noncommuting supports meet, if at all,
in a nonsingular line.
\end{proof}

\begin{theorem}[The six-dimensional string obstruction]\label{six}
A sequence of involutions in $\Omega(V)$ with the string commutations and
with noncommuting adjacent pairs generates a reducible group on $V$.
\end{theorem}
\begin{proof}
Let $U_i=\im(g_i-I)$. The sum of all these supports is invariant under all
generators, since $g_i(v)-v\in U_i$ for every $v$. For at most two generators
its dimension is at most four. For three it is at most five, because $U_0$
meets $U_2$. For four, both $U_0+U_2$ and $U_1+U_3$ have dimension at most
three and they intersect, since $U_0$ meets $U_3$. Their sum is again proper.

Assume at least five generators. If $U_0,U_1$ are disjoint, each $U_j$
for $j\ge3$ meets both and lies in their four-dimensional sum. Since
$U_2$ meets $U_0$, the total support has dimension at most five.
Otherwise $U_0\cap U_1=\gen b$ with $Q(b)\ne0$ by
Lemma~\ref{support}. Put $W=U_0+U_1\subset b^\perp$.

If $U_1,U_2$ are disjoint, put $W'=U_1+U_2$.
All $U_j$, $j\ge4$, lie in $W'$, and $U_0$ meets $U_2$.
If $U_0\subset W'$, the plane $U_3$ meets $U_1$ and adds at most one
dimension. Otherwise $W'+U_0$ has dimension five. The plane $U_3$ meets
both $U_1$ and $U_0$, so lies in $W'+U_0$ unless those two intersections
are the same line. In the latter case that line lies in $U_1\cap U_0$,
which is disjoint from $U_2$; since $U_0\cap W'$ already contains the
line $U_0\cap U_2$, two distinct lines would force $U_0\subset W'$,
a contradiction. Thus the total support is again proper.

We may therefore suppose $U_1\cap U_2=\gen c$ with $Q(c)\ne0$.
The plane $U_2$ also meets $U_0$. If it contains $b$, then $U_2\subset b^\perp$;
otherwise its two intersections span it and $U_2\subset W\subset b^\perp$.
Thus $g_0,g_1,g_2$ fix $b$. Each later $U_j$ meets both $U_0,U_1$;
it either contains $b$ or lies in $W$. In both cases it is perpendicular
to $b$, so $g_j$ fixes $b$. The whole generated group has a fixed line.
\end{proof}

\begin{proposition}\label{four}
For every even prime power $q$, neither $\operatorname{PSL}_4(q)$ nor
$\operatorname{PSU}_4(q)$ admits a string C-group marking of any rank.
\end{proposition}
\begin{proof}
Retain the standard Klein identifications
$\operatorname{PSL}_4(q)\cong\Omega^+_6(q)$ and
$\operatorname{PSU}_4(q)\cong\Omega^-_6(q)$, with irreducible natural
six-dimensional actions. In these even-characteristic parameters both are
nonabelian simple. A commuting adjacent pair in a string C-group would split
it as a direct product of two nontrivial subgroups, contrary to simplicity.
Theorem~\ref{six} now excludes a generating marking. This proves the minus-type
transfer explicitly, including $q=2$; it does not identify the simple derived
group with the full orthogonal group.
\end{proof}

\begin{thebibliography}{9}
\bibitem{LV} D. Leemans and A. Vandenschrick, \emph{On simple groups that admit
a string C-group representation}, locally bound six-page author manuscript,
\url{https://dipot.ulb.ac.be/dspace/bitstream/2013/397490/3/simpleGroups6.pdf}.
\bibitem{BV} P. A. Brooksbank and D. A. Vicinsky, \emph{Three-dimensional
classical groups acting on polytopes}, Discrete Comput. Geom. \textbf{44}
(2010), 654--659. Six-page author version, 28 June 2009,
\url{https://www.unix.bucknell.edu/~theburg/projects/reports/vicinsky.pdf}.
\bibitem{BL} P. A. Brooksbank and D. Leemans, \emph{Polytopes of large rank
for $\operatorname{PSL}(4,q)$}, J. Algebra \textbf{452} (2016), 390--400.
The 26 November 2015 author manuscript, Section~4, was read in the author's
ResearchGate full-text representation. It is the fixed-catalogue L048 cited
as a proof prerequisite, not a new recovery or source count.
\bibitem{Wilson} R. A. Wilson, \emph{The Finite Simple Groups}, Springer, 2009,
Section~3.11.2 for the unitary Klein correspondence.
\end{thebibliography}
\end{document}
